\chapter{Introduction}\label{ch:introduction}
\thispagestyle{pagestyle}

\section{Motivation}
Quantum computing is moving from proof-of-principle algorithms toward applications in signal and image processing~\cite{Ruan2021OpportunitiesChallenges,Wang2022ReviewQIP,Farooq2025SystematicReview}. Quantum image representations encode pixel data in amplitudes or basis states of a small register, enabling quantum versions of filtering, segmentation, and edge detection~\cite{Le2011FRQI,Yao2017QHED}. On current noisy intermediate-scale quantum (NISQ) hardware, however, circuit depth, two-qubit gate count, and readout errors dominate feasibility long before asymptotic speedups matter~\cite{Geng2022HybridNISQEdge,Geng2023ImprovedFRQI}.

This thesis studies a concrete pipeline: \emph{flexible representation of quantum images} (FRQI) encoding~\cite{Le2011FRQI} followed by the quantum Hadamard edge detector (QHED) baseline~\cite{Yao2017QHED,QiskitTextbookQuantumEdgeDetection}. The scientific question is not a new edge operator, but whether a \textbf{structurally cheaper FRQI state-preparation circuit}---obtained by changing multi-controlled-$X$ synthesis---changes transpiled resources and robustness under noise compared with a textbook naive decomposition.

\section{Baseline and research gap}
FRQI stores each grayscale intensity in a rotation angle on a dedicated color qubit, while a position register of $m = 2\log_2 N$ qubits addresses an $N\times N$ image~\cite{Le2011FRQI,Iliyasu2013FRQIroadmap}. Implementations that loop over pixels and apply multi-controlled rotations naively inherit worst-case multi-controlled gate costs; recent synthesis literature shows that ancilla-assisted decompositions can reduce CX depth for multi-controlled targets~\cite{Vale2023MCSU2,Azevedo2024LowDepthMCG,Arsoski2024MCXQFT,Zindorf2024EfficientMCG}. Separately, readout errors on the color qubit bias estimated intensities; calibration-matrix inversion is a standard mitigation primitive on superconducting backends, and we include a controlled simulator slice to illustrate the workflow~\cite{Mu2025GuidedImageFiltering}.

Prior quantum edge-detection work often assumes ideal state preparation or focuses on NEQR-based Sobel variants~\cite{Zhang2013NEQR,Zhang2015QSobel,Fan2019SobelNEQR,Chetia2021ImprovedSobelNEQR}. Few studies pair \textbf{paired} naive versus ancilla-assisted FRQI preparation under identical noise grids while holding the downstream QHED stage fixed.

\section{Contributions}
This dissertation makes three contributions, all supported by reproducible scripts and CSV manifests in the project repository:
\begin{enumerate}
  \item A \textbf{structural comparison} of naive (\texttt{noancilla}) versus v-chain (\texttt{v-chain}) multi-controlled FRQI preparation, including transpiled qubit count, depth, and CX for $4\times 4$, $8\times 8$, and $16\times 16$ test images, with honest scaling estimates where full naive transpilation is infeasible.
  \item A \textbf{noisy-simulator robustness study} that reconstructs FRQI states under depolarizing noise, applies QHED, and reports PSNR, SSIM, fidelity, and edge-map metrics for naive and v-chain layouts on matched noise grids.
  \item A \textbf{readout mitigation proof-of-concept} that estimates a $2\times 2$ confusion matrix for the color qubit, inverts it with regularization, and compares raw versus mitigated $p_1$ statistics across random seeds on a $4\times 4$ slice.
\end{enumerate}

\section{Scope and limitations}
All experiments in this submission use \textbf{simulators only} (Qiskit Aer statevector or density-matrix backends, plus optional mock NISQ noise wrappers). No IBM Quantum hardware runs are included. Image sizes are limited to $4\times 4$, $8\times 8$, and $16\times 16$ grayscale test patterns. Density-matrix simulation of the full $16\times 16$ v-chain register (16~qubits) is often skipped to respect memory limits; the Results chapter states where v-chain rows are missing. Statistical tests on paired naive--v-chain differences are reported as \textbf{exploratory} and should not be read as definitive superiority claims.

\section{Thesis outline}
\Cref{ch:related} surveys quantum image representations, edge-detection literature, NISQ constraints, multi-controlled synthesis, and mitigation. \Cref{ch:methods} describes methods, experimental setup, and results. \Cref{ch:conclusions} summarizes contributions, limitations, and future work including hardware validation and zero-noise extrapolation.
